\chapter{Condensed Matter}
\label{chap:condensed-matter}

Condensed matter is where the instrument is at its most modest, and that modesty is itself a kind of honesty.
Here, for the first time, the physics outruns the finite model so far that the model cannot pretend to derive
it: the band structure of a crystal, the pairing of electrons, the expulsion of a field are collective,
many-body facts the instrument cannot read off its counts the way it read a trajectory or a balance. And so
the instrument does the honest thing --- it names the effect and claims no more. The name is ours to give; the
collective fact it points at is the world's, outrunning here anything the model can derive. This is the fourth
and last of its honest floors, the name-only floor, and the instrument marks its debut as carefully as it
marked the wall and the analogy before it. Where a finite count is a balance the instrument secures, a smooth
shadow an analogy it names, and a no-go a wall it proves, the name-only floor is barer than all three: the
phenomenon is named, the physics is written out in full as physics, and the finite model claims only the label
beneath it --- the label ours, the collective fact beneath it the world's. Three effects mark the new floor
and, with it, complete the instrument's set of four: the semiconductor's bandgap, the superconductor's paired
state, and the Meissner expulsion of flux.

\section{The Semiconductor Effect: bandgap conduction}
\label{se:semiconductor}

A solid's electrons fill its available states up to the Fermi level, and whether the solid conducts depends
on what sits at that level. In a metal the level falls in the middle of a band and electrons move freely; in
an insulator a forbidden bandgap $E_g$ separates a filled valence band below from an empty conduction band
above, and no available variation crosses it. A semiconductor is the intermediate case --- a gap small enough
that thermal energy can lift a few electrons across it. The carrier density is thermally activated,
\begin{equation}
  n \propto e^{-E_g/2k_B T},
\end{equation}
climbing as the temperature rises and carriers are promoted over the gap, while the Fermi level sits within
the gap, partitioning the filled bands below from the empty bands above. The occupation of the states follows
the Fermi--Dirac distribution,
\begin{equation}
  f(E) = \frac{1}{e^{\beta(E-\mu)} + 1},
\end{equation}
the probability that a state of energy $E$ is filled at temperature $T$, near one for states well below the
Fermi level $\mu$ and near zero well above it, with $\mu$ the half-occupancy line where $f = \tfrac12$. The
carrier density above the gap is that distribution's exponential tail, which is why $n$ climbs as $e^{-E_g/2k_B
T}$. On this small fact --- a gap that thermal energy can bridge, a conduction a voltage or a dopant can switch --- the entire electronic age is
built. Bardeen and Brattain made the first transistor from a sliver of germanium in 1947, and every chip since
is the semiconductor's bandgap, switched.

Here the instrument meets the barest of its honest floors, the name-only floor. The activation law above is
the physics, offered as physics; the finite model claims nothing but the label. It records that the effect is
named --- ``the semiconductor effect'' --- and models none of the band dynamics that produce it. The name is
ours to give; the band-structure fact it names --- which no counting of ours derives --- is the world's. The
reading is that the semiconductor is named, the activation law $n \propto e^{-E_g/2k_B T}$ is the physics, and
the whole of the finite model's claim is the name beneath it. This is not a retreat but an honesty: the
instrument does not pretend to a finite model of band structure it does not have.

\section{The Superconducting Effect: the paired state}
\label{se:superconducting}

Below a critical temperature, the electrons of certain metals do something no classical picture allows: they
pair. Two electrons $(e_i, e_j)$ that would each, alone, scatter off the lattice and resist a current bind
into a single coherent pair, and an
energy gap $\Delta$ opens that protects the pairs from the scattering that would slow them --- $\Delta$ the
energy it costs to break a single pair, and $2\Delta$ the gap a spectroscope reads in the excitation spectrum,
since breaking a pair frees two electrons. The result is
superconduction: current flows without loss, at exactly zero electrical resistance, for as long as the metal
is kept cold. Bardeen, Cooper, and Schrieffer gave the microscopic account in 1957 --- the pairs now bear
Cooper's name --- and the effect is visible to the eye: a magnet set above a cooled superconductor hangs in
the air, held up by a field the superconductor will not let in.

The honest floor is again name-only, and here the discipline takes a particular care. The pairing and the gap
$\Delta$ are the physics, written as physics; the finite model stands only as a label beneath them, naming the
effect and claiming nothing of its mechanism. The name is ours to give; the coherent paired state the world
holds together --- which no counting of ours derives --- is the world's. The pair algebra that a finite
construction might build around such a state --- the bookkeeping of how paired carriers compose --- is not this
section's subject, and is not developed here. Here the superconductor is named, and the physics is left as
physics. The reading is that the superconductor is named, the gap and the paired state are physics, and the
finite model is the label and nothing more.

\section{The Meissner Effect: flux expelled}
\label{se:meissner}

A superconductor does not merely fail to admit a magnetic field; it expels one. Cool a superconductor below
its transition in the presence of a field, and the field is pushed out of the interior --- not frozen at
whatever value it held, but actively expelled, so that inside the material it vanishes,
\begin{equation}
  \mathbf{B} = 0 \quad\text{inside } \Omega, \qquad F_{\mu\nu}\big|_{\Omega} = 0,
\end{equation}
the field tensor zero throughout the interior. The superconductor is a perfect diamagnet, screening its
interior with surface currents that exactly cancel the applied field within. This is the expulsion that
levitates the magnet: the superconductor will not hold the field inside it, so it pushes back on the magnet
that would impose one, and the magnet floats. The Meissner effect is the sharper statement of
superconductivity --- not merely that current flows without resistance, but that the field is driven out, a
true thermodynamic phase and not just a perfect conductor.

The honest floor is name-only. The expulsion $\mathbf{B} = 0$ is the physics, stated as physics; the finite
model is the label beneath it, naming the effect and claiming no model of the screening currents that enforce
it. The name is ours to give; the flux the world drives out --- which no counting of ours derives --- is the
world's. The reading is that the Meissner effect is named, the vanishing of the interior field is physics, and
the finite model is the name --- the third effect, and the last, to stand on the barest of the instrument's
four floors.

It is worth saying why the name-only floor appears here and not earlier. The instrument counted atoms in chemistry
and balanced reactions in whole numbers, and it could do that because chemistry's facts are, at
bottom, the facts of individual atoms. But a semiconductor's bandgap, a superconductor's pairing, a perfect
diamagnet's expulsion of a field are not facts about any one particle --- they are emergent, the collective
behavior of countless particles acting together, and no count of individual events derives them. This is the
limit of reduction, and the instrument meets it honestly. It does not pretend that its ledger of single
events explains the band structure or the pairing; it names the emergent effect, writes its effective physics
in full, and stands on the barest floor. The name-only floor is the instrument's honest answer to emergence:
where the whole is more than the count of its parts, name the whole and leave the physics as physics.

\bigskip

With the name-only floor, the instrument's set of honest floors is complete, and all four are now in its
hands: a finite count obligation, a balance it secures; a smooth-shadow analogy, a continuum it names as a
shadow; a no-go wall, a refinement it proves impossible; and a name-only label, a phenomenon it names and
leaves to physics. For every effect still to come, the instrument will stand on exactly one of
these four, and say which. And with the solids, Part III closes. Heat and matter, read off finite records,
have been counts throughout: entropy is a count of distinctions, a domain wall is a no-go, a temperature is
the stable shape of accumulated error, a pressure is the density of reconciliation at a wall, and a solid's
order is a name the count settles into. Part IV turns next to the quantum, where the count becomes an amplitude and the amplitude, squared,
becomes a probability. The doorway Brownian
motion opened, diffusion run a quarter-turn into imaginary time, waits there.

\bigskip
\begin{center}
\rule{0.4\textwidth}{0.4pt}
\end{center}
\bigskip

\noindent Some evidence the record can only name. A solid's order --- the paired state, the expelled field, the
gap that switches a current --- is the work of countless parts acting as one, and no counting of the parts hands
it over. So the honest instrument enters the effect by name, writes its physics in full, and claims nothing
beneath the label. The name is ours to give; the collective fact it points at, which no tally of ours derives,
is the world's. An honest account admits plainly what it can only call by name.
